\documentclass[../main.tex]{subfiles}

\begin{document}

\chapter{System Design}\label{ch:system-design}

\section{Requirements for the System}\label{sec:requirements}
What comes after the research questions and the related work:
- Four requirements: multiple independent sources, because a single source does not allow for a comparison of actual versus target figures. 
- passive monitoring without acting on alerts (Looking-Glass Principle)
- no source code editing in Shibboleth (work as a plugin), so that it remains a deployment artifact
- Storage outside the IdP's trust boundary, with implications for transport and pseudonymization.
- Store saved data on a differente trust boundary of a IdP

\section{Overview}\label{sec:overview}
- The four layers and their coordiation / how they work together
- Architecture diagram

\section{Data Aggregation Layer}\label{sec:dataaggregation}
- Through what kinds of channels does an application actually disclose information—five types.
- Two can be ruled out because they cannot be linked to a specific login. This answers RQ1.
- The three data sources (audit, config, interceptor), with the questions: 
    - What does it see, and why does it see it? 
    - Explain the different interceptor points when explaining about the intercepor
        - what would I see when I would intercept it earlier
        - or later in the hook
- the dynamic field adaption
- the admin decides the audit format / what is being logged
- Then, in sections 4.5 and 4.6, simply address the implications for rules and correlation.

\section{Transport Layer and Trust Boundary}\label{sec:trust}
- Why do I store the data outside of the trust boundary
- HMAC, Pseudonymization at the source
- And the decision to use asynchronous transmission, with the rationale that the analysis must not slow down the monitored service

\section{Correlation Layer}\label{sec:correlation}
- Why is the correlation necessary?
- Two independent accounts of the same event.
- Explain the idea of the Fallback-Chain
- Why a chain and not a strategy
    - the admin decides what is being logged
- without implementation details

\section{Detection Layer}\label{sec:detection}
- How is a rule build
- the four rule engines and why four
    - the differ in the input
- And the rollback in the event of missing fields, as a consequence of 4.3.

\section{Design Decision and Rejected Alternatives}\label{sec:design}
- why no automatic response
- no intervention in the logging chain
- no machine learning.
- For each decision: What would have been the alternative, and why not?

\end{document}